\documentclass[11pt,twoside]{article}

\usepackage[margin=1in]{geometry}
\usepackage{amsmath,amssymb,amsthm,mathtools}
\usepackage[inline,shortlabels]{enumitem}
\usepackage{xcolor}
\usepackage{../misc_files/course-definition-boxes}
\usepackage{microtype}

% Revision date: maintained by scripts/build_lecture_notes.py. Keep this command.
\newcommand{\lecturerevised}{2026-11-24 22:19 UTC-06:00}
\usepackage{eso-pic}
\AddToShipoutPictureBG{%
  \AtPageLowerLeft{%
    \put(\LenToUnit{0.5\paperwidth},\LenToUnit{0.5in}){%
      \makebox(0,0){\normalfont\footnotesize\color{black!60}Last revised: \lecturerevised}%
    }%
  }%
}

\usepackage[square,authoryear]{natbib}
\usepackage[colorlinks=true,linkcolor=blue!60!black,
  urlcolor=blue!60!black,citecolor=blue!60!black]{hyperref}
\usepackage[nameinlink,capitalize]{cleveref}
\crefname{exercise}{Exercise}{Exercises}
\Crefname{exercise}{Exercise}{Exercises}
\usepackage{../exercises/course-exercises}

\usepackage{enotez}
\DeclareInstance{enotez-list}{lecture}{list}{format=\normalsize,number={#1.}}
\setenotez{list-name=Endnotes,list-style=lecture,backref=true}
\crefname{endnote}{Endnote}{Endnotes}
\Crefname{endnote}{Endnote}{Endnotes}

\definecolor{courseblue}{HTML}{17365D}
\newcommand{\E}{\mathbb E}
\newcommand{\I}{\mathbb I}
\renewcommand{\P}{\mathbb P}
\newcommand{\R}{\mathbb R}
\newcommand{\Law}{\mathcal L}
\newcommand{\KL}{D}
\newcommand{\TV}{\mathrm{TV}}
\makeatletter
\@for\@letter:=A,B,C,D,E,F,G,H,I,J,K,L,M,N,O,P,Q,R,S,T,U,V,W,X,Y,Z\do{%
  \expandafter\edef\csname c\@letter\endcsname{\noexpand\mathcal{\@letter}}}
\makeatother
\newcommand{\lecturenumber}{22}
\newcommand{\suppmark}{\texorpdfstring{\textsuperscript{*}}{*}}
\newcommand{\xor}{\mathbin{\oplus}}
\DeclareMathOperator*{\argmin}{arg\,min}

\newtheorem{theorem}{Theorem}[section]
\newtheorem{proposition}[theorem]{Proposition}
\newtheorem{lemma}[theorem]{Lemma}
\newtheorem{corollary}[theorem]{Corollary}
\theoremstyle{definition}
\newtheorem{definition}[theorem]{Definition}
\newtheorem{example}[theorem]{Example}
\theoremstyle{remark}
\newtheorem{remark}[theorem]{Remark}

\renewcommand{\thepage}{\lecturenumber--\arabic{page}}
\renewcommand{\thesection}{\lecturenumber.\arabic{section}}
\renewcommand{\theequation}{\lecturenumber.\arabic{equation}}

\pagestyle{myheadings}
\markboth{Lecture \lecturenumber: Statistical and exact learning}
{Lecture \lecturenumber: Statistical and exact learning}

% Course-wide notation convention:
% Consult notation-reference.md and the source definitions it indexes before
% introducing symbols; repeat the comparison on the completed draft.
% See learning-notation-conventions.tex: \ell is one-observation loss;
% L is population risk and L_n is empirical average loss. Keep the standard
% term empirical risk minimization (ERM). Parameterization does not change
% this convention. State local exceptions and all normalization factors.
% Every vector is a column vector. If a matrix stores one vector per position,
% row i contains the transpose of that vector. Write every row-shaped vector
% explicitly as the transpose of a column vector.
% Use n for the number of independent training examples. Use sentence case in
% optional theorem titles. Use example environments for worked examples that
% are referenced later. Reserve \boxed{} for literal boxes in diagrams.

\begin{document}
\raggedbottom

\begin{center}
\fcolorbox{courseblue}{white}{%
\begin{minipage}{0.92\textwidth}
\vspace{1mm}
\textbf{CMPUT 654: Mathematical Foundations of Modern AI Systems}
\hfill Fall 2026

\vspace{3mm}
\begin{center}
{\Large\bfseries Lecture \lecturenumber: Statistical and exact learning}

\vspace{1mm}
November 24, 2026
\end{center}

\vspace{2mm}
\textit{Lecturer: Csaba Szepesv\'ari}
\hfill
\textit{Notes prepared by: Csaba Szepesv\'ari}
\vspace{1mm}
\end{minipage}}
\end{center}

Suppose that a designer wants to teach one specified rule to a machine. The
designer may be able to label or verify individual examples without knowing
how to construct a predictor that applies the rule to every valid input. A
small average error then answers only part of the design question. The learned
predictor may still implement a different rule on inputs that are rare under
the training distribution. A two-rule argument first shows that passive data
can leave some target unidentified. Learner symmetry then converts that
class-level statement into a lower bound for the one intended rule. The same
symmetry is often desirable: a representation convention should not decide a
question that the data or task specification is meant to decide. The
principal theorem quantifies the resulting cost and identifies the information
that a teacher, query, architecture, or verifier must supply.

\noindent\textbf{Reading guide.}
The distinction between worst-case and application-specific guarantees from
\href{lecture02.pdf}{Lecture 2}, teaching and active learning from
\href{lecture06.pdf}{Lecture 6}, and maximum-margin implicit bias from
\href{lecture10.pdf}{Lecture 10} provide the main prerequisites. The
maximum-margin teaching proof is deferred to an appendix. An optional appendix
revisits the reverse-KL support calculation from
\href{lecture04.pdf}{Lecture 4}. The system-level discussion connects the
result to verification from
\href{lecture17.pdf}{Lecture 17} and selective queries from
\href{lecture18.pdf}{Lecture 18}.

\section{Teaching one intended rule}
\label{sec:one-rule}

The engineering objective in this lecture is one specified competence. The
designer knows which behavior is intended, at least well enough to label or
verify instances, and chooses a method intended to produce that behavior. The
designer does not require the same method to learn every rule in a broad
family.

Let \(\cX\) be a finite set of valid inputs and let
\(h^\star:\cX\to\{0,1\}\) be the intended rule. Fix an input law \(\rho\)
on \(\cX\), and let \(X\) have law \(\rho\). For a candidate rule
\(h:\cX\to\{0,1\}\), use zero--one loss
\[
 \ell(h,(x,y))=\I\{h(x)\ne y\}.
\]
The population risk under \(\rho\) is
\begin{equation}
 L_\rho(h)=\E[\ell(h,(X,h^\star(X)))]
 =\P(h(X)\ne h^\star(X)).
 \label{eq:population-error}
\end{equation}
The subscript remains visible because the argument will compare this one-law
average with correctness under other input laws.

\begin{definition}[Statistical and exact correctness]
The rule \(h\) has \emph{statistical error at most \(\varepsilon\)} under
\(\rho\) when \(L_\rho(h)\le\varepsilon\). The rule \(h\) is
\emph{exactly correct} for the specified task when
\begin{equation}
 h(x)=h^\star(x)\qquad\text{for every }x\in\cX.
 \label{eq:exact-rule}
\end{equation}
\end{definition}

Exact correctness changes the quantifier over inputs without introducing a
quantifier over target rules. The distinction is central: requiring
\cref{eq:exact-rule} for the one intended \(h^\star\) does not require a
learning procedure that succeeds for every map from \(\cX\) to
\(\{0,1\}\).

\begin{proposition}[Exactness as distributional robustness]
\label{prop:all-laws}
Let \(\cM(\cX)\) be the set of all probability laws on \(\cX\), and define
\[
 L_\nu(h)=\sum_{x\in\cX}\nu(x)\I\{h(x)\ne h^\star(x)\}
 \qquad\text{for }\nu\in\cM(\cX).
\]
Then
\begin{equation}
 \sup_{\nu\in\cM(\cX)}L_\nu(h)
 =\max_{x\in\cX}\I\{h(x)\ne h^\star(x)\}.
 \label{eq:all-laws}
\end{equation}
Consequently, \(h\) is exactly correct if and only if the left side of
\cref{eq:all-laws} is zero.
\end{proposition}
\begin{proof}
Every average of the indicator is at most its maximum, so the left side is at
most the right side. Conversely, let \(x_0\) attain the maximum and take
\(\nu=\delta_{x_0}\). Then
\(L_{\delta_{x_0}}(h)=\I\{h(x_0)\ne h^\star(x_0)\}\), which gives the reverse
inequality.
\end{proof}

The exact criterion is appropriate when valid inputs have a
distribution-independent specification of correctness. Examples include
executing an algorithm, applying formal rules, checking a formal proof, and
following a protocol. Statistical prediction remains appropriate when the
target is intrinsically stochastic or when no unique correct response has
been specified. A system may contain both kinds of component.

\endnote{Changing the KL direction does not reveal the intended response at an
unseen input. For a deterministic full-support target law, finite reverse KL
characterizes exact conditional support, but every imperfect candidate has
infinite divergence. Forward KL is finite and graded when the candidate assigns
positive mass to each correct response, but it averages log loss under the input
law. Replacing the target law by an empirical law only restricts the candidate
to observed atoms. A graded reverse-KL objective requires a softened target or
an evaluable cost, and either choice adds information or assumptions. The
optional comparison appears in \cref{app:kl-directions};
\cref{ex:l22-reverse-kl-exactness} develops the calculations.}

\section{What passive examples can distinguish}
\label{sec:pairwise}

The first lower bound concerns two possible targets. The result is deliberately
learner-independent: if the sample never enters the region where two rules
disagree, no procedure can infer from that sample which rule supplied the
labels.

Let \(\cA\) be a possibly randomized learner that maps a labeled dataset to a
binary rule. For a rule \(h:\cX\to\{0,1\}\), draw \(X_1,\ldots,X_n\)
independently with common law \(\rho\), and define the labeled sample
\begin{equation}
 D_h=((X_1,h(X_1)),\ldots,(X_n,h(X_n))).
 \label{eq:labeled-sample}
\end{equation}
The learner's output on \(D_h\) is the random rule \(H_{D_h}\); the randomness
comes from the sample and any private randomness used by the learner. The
pointwise exact-learning loss is
\begin{equation}
 \ell_{\mathrm{ex}}(g,h)=\I\{g\ne h\},
 \label{eq:exact-learning-loss}
\end{equation}
where inequality means inequality as functions on \(\cX\). Define the expected
exact-learning loss, or exact-recovery failure probability, by
\begin{equation}
 \phi_n(\cA,\rho,h)
 =\E[\ell_{\mathrm{ex}}(H_{D_h},h)]
 =\P(H_{D_h}\ne h).
 \label{eq:failure-probability}
\end{equation}
The letter \(\phi\) is only a mnemonic for failure; no standard convention
requires that notation.

\begin{lemma}[Pairwise indistinguishability]
\label{lem:pairwise}
Let \(h,h':\cX\to\{0,1\}\) be distinct, and write
\begin{equation}
 \alpha_\rho(h,h')=\P(h(X)\ne h'(X)).
 \label{eq:pair-disagreement}
\end{equation}
Every learner \(\cA\) satisfies
\begin{equation}
 \phi_n(\cA,\rho,h)+\phi_n(\cA,\rho,h')
 \ge (1-\alpha_\rho(h,h'))^n.
 \label{eq:pairwise-bound}
\end{equation}
\end{lemma}
\begin{proof}
Let
\[
 E=\bigcap_{i=1}^n\{h(X_i)=h'(X_i)\}.
\]
Independence gives \(\P(E)=(1-\alpha_\rho(h,h'))^n\). On \(E\), the two
labeled datasets \(D_h\) and \(D_{h'}\) are identical. Conditional on any
common dataset, one output rule cannot equal both distinct rules. Therefore
the two conditional failure probabilities sum to at least one. Averaging this
inequality on \(E\) gives \cref{eq:pairwise-bound}.
\end{proof}

The usual hypothesis-class consequence follows immediately. For a set
\(\cH\) of candidate rules, define its \emph{worst-target failure probability}
by
\begin{equation}
 \Phi_n(\cA,\rho,\cH)
 =\sup_{h\in\cH}\phi_n(\cA,\rho,h).
 \label{eq:class-failure}
\end{equation}

\begin{corollary}[A hypothesis class containing a close pair]
\label{cor:class-pair}
If distinct \(h,h'\) belong to \(\cH\), then every learner satisfies
\begin{equation}
 \Phi_n(\cA,\rho,\cH)
 \ge\frac12(1-\alpha_\rho(h,h'))^n.
 \label{eq:class-pair-bound}
\end{equation}
\end{corollary}
\begin{proof}
Each failure probability in \cref{eq:pairwise-bound} is at most
\(\Phi_n(\cA,\rho,\cH)\).
\end{proof}

The corollary says that at least one member of a close pair is hard. The
conclusion is useful when one procedure must work uniformly over \(\cH\), but
the engineering problem in this lecture has different quantifiers. The target
\(h^\star\) and the learner actually used are both fixed. The corollary allows
its hard target to depend on the learner.

A learner that always returns \(h^\star\) has
\begin{equation}
 \phi_n(\cA,\rho,h^\star)=0
 \qquad\text{for every }n,
 \label{eq:hard-coded-target}
\end{equation}
even though the same learner can fail with probability one on every other
target. The pairwise sum remains large, so neither the lemma nor the class
corollary rules out immediate success on the intended rule.

The constant learner would be an excellent solution if the designer could
implement the complete rule \(h^\star\). The premise of the teaching problem is
that the designer can label or verify individual instances but does not possess
that complete implementation. Target-dependent examples, features, and
architectures should remain available; a lower bound should not prohibit useful
target information.

Ideally, the analysis would determine the exact-learning difficulty of the
particular training procedure that will be used. A complete analysis of one
large training system can be intractable. The next argument therefore retains
one structural property of that procedure and proves a target-specific lower
bound for every learner with the property. The mathematical class now consists
of learners rather than target rules. The resulting statement is weaker than a
complete analysis of the particular learner because it uses only the retained
property. Learner symmetry is the property used here, and the symmetry must be
verified for the task and learner under study.

\section{Why and when to require learner symmetry}
\label{sec:symmetry}

Symmetry defines a class of learners through a property of their input--output
behavior. A flexible learning procedure can represent many rules, but
representability does not explain which rule the procedure selects. The role of
learner symmetry is to state which choices should be determined by data or by
an explicit task specification, rather than by features of the representation
or implementation that the task declares irrelevant.

\subsection{Why learning creates a role for symmetry}
\label{sec:why-symmetry}

The designer in this lecture knows the intended behavior well enough to label,
verify, or deliberately select individual instances. The designer nevertheless
lacks a complete executable predictor, or chooses not to encode one manually.
Learning delegates some choices to data. A suitable learner should therefore
avoid making those choices from information that the designer regards as
irrelevant.

For example, the numerical names \(0\) and \(1\) carry no semantic information
in an ordinary binary classification problem. Consistently exchanging the
names in the training data should exchange the predicted labels. Likewise,
the storage order of an i.i.d.\ batch should not affect a batch learner. For
object classification, translating an object within an image may be intended
to preserve its class. In each case, symmetry prevents a convention---a class
code, record order, or image origin---from deciding the learned behavior.

The same transformations can be invalid for a different task. Label exchange
is inappropriate when the labels have different error costs. Record order is
part of the evidence in a time series or conversation. Translation is not a
task symmetry when the target asks whether an object occurs on the left or the
right. A symmetry must therefore come from the semantics and information
structure of the task; visual or algebraic similarity alone is insufficient.

The batch-order transformation acts on example indices rather than on the
input space \(\cX\). The batch-order example illustrates the design principle
but is not an element of the input--label group used in the theorem below.
Label renamings and input translations do fit that formalism.

Choosing a symmetry is itself an inductive bias. The choice asserts that
alternatives related by the transformation deserve equal treatment before the
data or specification distinguish them. A learner may and should contain
other asymmetries supported by prior task knowledge. If the specification
states that equality counts as a positive comparison, encoding that convention
is legitimate. If the convention is meant to be learned from examples, then a
tie example or an equivalent source of information must communicate it.

Equivariance also differs from refusing to choose. A transformed dataset may
lead to a transformed output. The requirement is that the two changes be
consistent. Data break a symmetry whenever the observed dataset is not itself
unchanged or statistically indistinguishable under the transformation. The
lower bound will apply precisely when passive data fail to provide that
asymmetry.

Before treating a transformation \(g\) as a required learner symmetry, the
designer should verify four conditions:
\begin{enumerate}[leftmargin=*]
\item \(g\) changes the representation while preserving the semantic task
under the corresponding input and output relabeling;
\item the observation or data-collection protocol transforms consistently
under \(g\);
\item the loss and evaluation criterion give corresponding transformed
decisions the same value; and
\item available prior knowledge does not justify preferring one
orbit-related alternative.
\end{enumerate}
When a condition fails, the corresponding asymmetry can be useful task
information. When all four conditions hold, breaking the symmetry through an
implementation accident lets the representation decide a question that was
delegated to learning.

\subsection{Formal setup}
\label{sec:symmetry-formal}

The formal definitions record consistent transformations of inputs, labels,
data, and learned rules. Let \(\cY:=\{0,1\}\), and let \(G\) be a group. Each
\(g\in G\) induces two bijections,
\[
 g_{\cX}:\cX\to\cX
 \qquad\text{and}\qquad
 g_{\cY}:\cY\to\cY,
\]
and these maps form group actions on their respective spaces. We use
\(g=(g_{\cX},g_{\cY})\) as shorthand for the two components of the task
transformation.

The action on rules is not a separate choice. Regard a rule
\(h:\cX\to\cY\) as its graph
\[
 \operatorname{graph}(h)=\{(x,h(x)):x\in\cX\}
 \subseteq\cX\times\cY.
\]
Transform the relation elementwise:
\[
 g\operatorname{graph}(h)
 =\{(g_{\cX}x,g_{\cY}h(x)):x\in\cX\}.
\]
Because \(g_{\cX}\) is bijective, the transformed relation is the graph of a
unique function. Denote that function by \(gh\). The defining commutative
identity and its equivalent inverse form are
\begin{equation}
 (gh)(g_{\cX}x)=g_{\cY}(h(x)),
 \qquad
 (gh)(u)=g_{\cY}\bigl(h(g_{\cX}^{-1}u)\bigr).
 \label{eq:function-action}
\end{equation}

The same pair of maps acts pointwise on a labeled dataset
\(D=((x_1,y_1),\ldots,(x_n,y_n))\):
\begin{equation}
 gD=((g_{\cX}x_1,g_{\cY}y_1),\ldots,
      (g_{\cX}x_n,g_{\cY}y_n)).
 \label{eq:data-action}
\end{equation}

\begin{definition}[Invariant input law]
The law \(\rho\) is \emph{\(G\)-invariant} when
\begin{equation}
 \Law(g_{\cX}X)=\rho
 \qquad\text{for every }g\in G
 \label{eq:invariant-law}
\end{equation}
whenever \(\Law(X)=\rho\).
\end{definition}

The learner may use a random initialization, randomized minibatches, or other
private randomness. Represent the learner by a probability kernel
\(\cA(\,\cdot\mid D)\) on the set of binary rules: after observing \(D\), its
output \(H_D\) has conditional law \(\cA(\,\cdot\mid D)\).

\begin{definition}[Equivariant learner]
The learner \(\cA\) is \emph{\(G\)-equivariant} when, for every dataset
\(D\), transformation \(g\in G\), and set \(B\) of output rules,
\begin{equation}
 \cA(gB\mid gD)=\cA(B\mid D),
 \qquad gB=\{gh:h\in B\}.
 \label{eq:learner-equivariance}
\end{equation}
Equivalently, the output on \(gD\) has the same law as the transformed output
\(gH_D\) on \(D\). The random rule \(gH_D\) is obtained by applying the
deterministic rule action in \cref{eq:function-action} to each realization of
\(H_D\).
\end{definition}

The definition describes the complete learner, including the model class. A
symmetry claim must account for the architecture, initialization, objective,
optimizer, stopping rule, preprocessing, and conversion from scores to
discrete predictions. Gradient methods provide an important class of examples
because their equivariance can be checked directly.

\subsection{Why gradient descent produces symmetric learners}
\label{sec:gradient-symmetry}

Gradient descent preserves a data transformation when the model and objective
admit a compatible orthogonal realization in parameter space. The task
transformation \(g=(g_{\cX},g_{\cY})\) is fixed first. The parameterization
must then supply a suitable map; such a map need not exist.

Let the parameter space be \(\R^p\). Fix \(g\in G\), and suppose that an
orthogonal matrix \(S_g\in\R^{p\times p}\) describes a compatible parameter
transformation. Let \(f_\theta\) denote the predictor with parameter
\(\theta\), and let \(F_D(\theta)\) be the training objective on dataset
\(D\). The two identities to check are
\begin{align}
 f_{S_g\theta}(g_{\cX}x)&=g_{\cY}\bigl(f_\theta(x)\bigr),
 \label{eq:predictor-parameter-action}\\
 F_{gD}(S_g\theta)&=F_D(\theta).
 \label{eq:objective-invariance}
\end{align}
The label action in \cref{eq:predictor-parameter-action} acts on the
predictor's output. For a probabilistic binary predictor, a label flip maps
the probability assigned to label one to its complement.

\begin{proposition}[Gradient-descent equivariance]
\label{prop:gd-equivariance}
Assume \cref{eq:predictor-parameter-action,eq:objective-invariance}, and let
\(S_g\) be orthogonal. Consider the updates
\begin{equation}
 \theta_{t+1}(D)=\theta_t(D)-\eta_t\nabla F_D(\theta_t(D)),
 \qquad t=0,\ldots,T-1,
 \label{eq:gd-update}
\end{equation}
where the stepsizes do not depend on the representation of the dataset. If
\begin{equation}
 \theta_0(gD)=S_g\theta_0(D),
 \label{eq:coupled-initialization}
\end{equation}
then
\begin{equation}
 \theta_t(gD)=S_g\theta_t(D)
 \qquad\text{for every }t=0,\ldots,T.
 \label{eq:coupled-iterates}
\end{equation}
Consequently, the learned predictor is \(G\)-equivariant. If the
initialization is sampled independently of \(D\) and its law is invariant
under every \(S_g\), the learner is equivariant as a probability kernel.
\end{proposition}
\begin{proof}
Differentiate \cref{eq:objective-invariance} with respect to \(\theta\).
Orthogonality gives
\begin{equation}
 \nabla F_{gD}(S_g\theta)=S_g\nabla F_D(\theta).
 \label{eq:gradient-transformation}
\end{equation}
Suppose that \(\theta_t(gD)=S_g\theta_t(D)\). Substituting
\cref{eq:gradient-transformation} into the update for \(gD\) gives
\[
 \theta_{t+1}(gD)
 =S_g\theta_t(D)-\eta_t S_g\nabla F_D(\theta_t(D))
 =S_g\theta_{t+1}(D).
\]
Induction from \cref{eq:coupled-initialization} proves
\cref{eq:coupled-iterates}. The prediction identity
\cref{eq:predictor-parameter-action} then proves equivariance of the output.
For an invariant random initialization, couple the two runs using
\(S_g\theta_0\) in the transformed run. The coupled initial parameters have
the required marginal laws, so the same argument proves equality in law.
\end{proof}

The proposition explains why gradient learners are central to the symmetry
lower bound. Once the model, loss, initialization, and update rule respect a
transformation, gradient descent does not introduce a preference between the
transformed problems. The preference must come from the data or from an
explicit asymmetric component of the training setup. The resulting neutrality
is desirable when the transformation represents an arbitrary convention, and
the same neutrality creates the exact-learning obstruction when passive data
rarely reveal which orbit element is intended.

The same proof applies to stochastic gradient descent when each per-example
loss satisfies the transformation identity and the two runs use the same
sequence of example indices. Random reshuffling and minibatch sampling then
preserve equivariance in law. Momentum, Euclidean gradient clipping, and
isotropic weight decay also preserve orthogonal transformations. Coordinatewise
adaptive methods, including Adam, preserve signed coordinate permutations but
need not preserve arbitrary rotations. A data-dependent stopping rule preserves
the symmetry only when its criterion has the corresponding invariance.

\subsection{Gradient-learner examples and nonexamples}
\label{sec:symmetry-examples}

The following examples distinguish transformations respected by the complete
training procedure from parameter redundancies and task transformations that
should not be imposed as symmetries.

\begin{example}[Coordinate permutations in a linear model]
Let \(P\) be a permutation matrix and let
\(z_{w,b}(x)=\langle w,x\rangle+b\). Replacing the inputs by \(Px\) and the
weight vector by \(Pw\) leaves the score unchanged:
\[
 z_{Pw,b}(Px)=z_{w,b}(x).
\]
An empirical objective built from these scores and an isotropic regularizer
therefore satisfies \cref{eq:objective-invariance}. An exchangeable
initialization law and gradient descent produce a learner that is equivariant
to input-coordinate permutations. Requiring this symmetry is justified only
when the coordinate names carry no task information. Known feature roles,
coordinate-specific units, or coordinate-specific regularization can supply a
legitimate asymmetry.
\end{example}

\begin{example}[Coordinate and label symmetry of a multilayer perceptron]
Consider a fully connected network with first-layer matrix \(W_1\), scalar
output score \(z_\theta(x)\), and binary cross-entropy loss
\begin{equation}
 \ell_{\mathrm{ce}}(z,y)=\log(1+\exp(z))-yz.
 \label{eq:binary-cross-entropy}
\end{equation}
For an input permutation \(P\), replace \(W_1\) by \(W_1P^\top\). The
transformed network satisfies
\[
 z_{S_P\theta}(Px)=z_\theta(x).
\]
If the first-layer columns have an exchangeable joint initialization law, the
resulting gradient learner has input-coordinate symmetry.

For a binary-label flip, negate the weights and bias in the scalar output
layer. The score changes from \(z\) to \(-z\), and
\begin{equation}
 \ell_{\mathrm{ce}}(-z,1-y)=\ell_{\mathrm{ce}}(z,y).
 \label{eq:cross-entropy-label-flip}
\end{equation}
A final-layer initialization whose law is invariant under sign reversal
therefore gives label symmetry. Standard independent, centered, symmetric
initializations satisfy both requirements. Gradient descent, minibatch SGD,
momentum, and Adam preserve the coordinate permutations and sign changes in
these two actions. Thus a standard gradient-trained multilayer perceptron can
have exactly the block-swap and label-flip symmetry used in
\cref{sec:comparison}. Any regularizer must also remain invariant under the
same parameter actions.

Label symmetry is appropriate when zero and one are interchangeable names.
An asymmetric evaluation cost, a distinguished safety label, or a
task-specific output convention can invalidate the symmetry even though the
network architecture still permits the parameter transformation.

The preceding statement applies directly to predicted probabilities. A hard
classifier also needs an equivariant treatment of a zero score. Random
tie-breaking with equal probabilities is label symmetric; the deterministic
rule \(\I\{z\ge0\}\) is not label symmetric at \(z=0\).
\end{example}

\begin{example}[Transformers]
A standard transformer is not equivariant to arbitrary permutations of
sequence positions: the causal mask and positional representation distinguish
positions. A transformer can, however, be equivariant to a consistent
renaming of vocabulary symbols when the renaming permutes the input-embedding
rows, output coordinates, and any tied embedding parameters. Exchangeable
initialization and a tokenwise cross-entropy objective preserve that action.
Fixed special tokens, a tokenizer with distinguished structure, or
token-specific preprocessing restrict the admissible renamings. A transformer
with a symmetric binary classification head can also have the label-flip
symmetry from the preceding example.
\end{example}

Two superficially similar symmetries do not yield a lower bound by themselves.
Permuting hidden units can leave a neural network's input--output function
unchanged. Such a parameter redundancy does not produce a distinct rule
\(gh^\star\). Conversely, an architecture may force every represented rule to
be invariant under a chosen input transformation. The transformation then
belongs to the target's stabilizer whenever the target is representable. The
target-specific lower bound uses a symmetry of the complete learner that maps
the intended rule to a distinct output rule.

The following audit is sufficient for applying the theorem to a gradient
learner:
\begin{enumerate}[leftmargin=*]
\item specify the action on inputs, labels, and output rules;
\item construct the corresponding parameter map \(S_g\) and verify
\cref{eq:predictor-parameter-action};
\item verify the objective identity \cref{eq:objective-invariance}, including
regularization;
\item check invariance of the initialization law and every optimizer state;
\item check minibatch sampling, stopping, preprocessing, and discrete
tie-breaking; and
\item exhibit a transformation for which \(gh^\star\ne h^\star\).
\end{enumerate}

For the fixed target \(h^\star\), define its \emph{orbit}
\begin{equation}
 \cO_G(h^\star)=\{gh^\star:g\in G\}.
 \label{eq:target-orbit}
\end{equation}
The orbit is not a collection of tasks that the designer wants to solve. The
orbit records alternatives that the learner treats as transformed copies of
the intended task. If the passive data law is also invariant, the learner has
equal failure probabilities across these alternatives.

\section{A target-specific symmetry lower bound}
\label{sec:lower-bound}

The main result now combines the two preceding developments. Pairwise
indistinguishability says that at least one of two targets must be hard.
Equivariance makes the orbit-related failure probabilities equal, so the lower
bound applies to the particular target that the designer intends to teach.

For \(g\in G\), define the \emph{target--orbit disagreement}
\begin{equation}
 \alpha_g(h^\star)
 =\P(h^\star(X)\ne(gh^\star)(X)).
 \label{eq:orbit-disagreement}
\end{equation}
Assume that the orbit contains a rule distinct from \(h^\star\), and define the
\emph{target-specific disagreement}
\begin{equation}
 \Delta_{\rho,G}(h^\star)
 =\inf_{g:\,gh^\star\ne h^\star}\alpha_g(h^\star).
 \label{eq:target-disagreement}
\end{equation}
The restriction \(gh^\star\ne h^\star\) removes transformations in the
target's stabilizer: those transformations do not create an alternative rule.

\begin{theorem}[Target-specific symmetry lower bound]
\label{thm:target-symmetry}
Suppose that \(\rho\) is \(G\)-invariant and \(\cA\) is
\(G\)-equivariant. For the fixed target \(h^\star\), every
\(g\in G\) with \(gh^\star\ne h^\star\) satisfies
\begin{equation}
 \phi_n(\cA,\rho,h^\star)
 \ge \frac12(1-\alpha_g(h^\star))^n.
 \label{eq:target-symmetry-bound}
\end{equation}
Consequently,
\begin{equation}
 \phi_n(\cA,\rho,h^\star)
 \ge \frac12(1-\Delta_{\rho,G}(h^\star))^n.
 \label{eq:target-symmetry-delta}
\end{equation}
If \(\Delta_{\rho,G}(h^\star)>0\) and
\(n<1/(2\Delta_{\rho,G}(h^\star))\), then
\begin{equation}
 \phi_n(\cA,\rho,h^\star)>\frac14.
 \label{eq:quarter-lower-bound}
\end{equation}
\end{theorem}
\begin{proof}
Fix \(g\) such that \(gh^\star\ne h^\star\). Transforming
\(D_{h^\star}\) by \(g\) gives a sample whose inputs are
\(g_{\cX}X_1,\ldots,g_{\cX}X_n\) and whose labels satisfy
\[
 g_{\cY}(h^\star(X_i))=(gh^\star)(g_{\cX}X_i)
\]
by \cref{eq:function-action}. Invariance of \(\rho\) therefore implies
\begin{equation}
 \Law(gD_{h^\star})=\Law(D_{gh^\star}).
 \label{eq:transformed-sample-law}
\end{equation}
Equivariance of the learner gives
\[
 \phi_n(\cA,\rho,gh^\star)
 =\phi_n(\cA,\rho,h^\star).
\]
Apply \cref{lem:pairwise} to \(h^\star\) and \(gh^\star\). The two failure
probabilities are equal, so each is at least half the right side of
\cref{eq:pairwise-bound}. The resulting inequality is
\cref{eq:target-symmetry-bound}.
Taking transformations with disagreement approaching the infimum gives
\cref{eq:target-symmetry-delta}.

Finally, \((1-z)^n\ge1-nz\) for \(z\in[0,1]\). If
\(n<1/(2\Delta_{\rho,G}(h^\star))\), then
\[
 \frac12(1-\Delta_{\rho,G}(h^\star))^n
 \ge\frac12(1-n\Delta_{\rho,G}(h^\star))
 >\frac14.
\]
\end{proof}

The theorem concerns one fixed target, one fixed learner, and one fixed data
law. The orbit in the proof is induced by the learner's symmetry. The theorem
does not require the designer to want one procedure that learns every element
of the orbit. Symmetry enters because the complete learner treats the orbit
elements as transformed versions of the same learning problem.

The proof also separates the roles of missing data and learner neutrality. The
pairwise lemma gives only
\[
 \phi_n(\cA,\rho,h^\star)
 +\phi_n(\cA,\rho,gh^\star)
 \ge(1-\alpha_g(h^\star))^n.
\]
A learner can satisfy this inequality by favoring \(h^\star\) before seeing
any data. Equivariance and input-law invariance give the additional identity
\[
 \phi_n(\cA,\rho,h^\star)
 =\phi_n(\cA,\rho,gh^\star),
\]
which divides the indistinguishability cost equally between the two targets.
The stronger conclusion is therefore a consequence of both an information
deficit in the sample and a neutrality property of the learner.

The neutrality property can be desirable. If \(g\) changes only a class code,
coordinate convention, or another representation choice, favoring one orbit
element would make that convention influence the learned rule. Removing the
symmetry arbitrarily builds in an answer to the unresolved choice. On a sample
that misses the disagreement region, an output chosen to equal \(h^\star\)
cannot also equal \(gh^\star\). Such a preference is appropriate when prior
task knowledge supports it; the preference should then be stated as part of
the learner or specification.

The theorem therefore gives a diagnostic and a design prescription. A
target-specific architecture can supply a justified asymmetry. A target-aware
curriculum can make the data law noninvariant. A membership query can
deliberately enter a disagreement region, and a verifier can search for one.
Each intervention states where information about the intended rule enters.
The desirable symmetry may remain in the learner while the teaching protocol
provides the evidence that breaks it.

\section{Teaching the comparison rule}
\label{sec:comparison}

The comparison between \(\ge\) and \(>\) gives an explicit fixed target for
which a coordinate-and-label symmetry has an exponentially small disagreement region. The
example separates label balance and full support from the information needed
for exact recovery.

Fix \(m\ge1\) and let \(\cX=\{0,1\}^{2m}\). For
\(x=(x_1,\ldots,x_{2m})\), define the two encoded integers
\begin{equation}
 L(x)=\sum_{j=1}^m2^{m-j}x_j,
 \qquad
 R(x)=\sum_{j=1}^m2^{m-j}x_{m+j}.
 \label{eq:encoded-integers}
\end{equation}
The intended rule is
\begin{equation}
 h_{\ge}(x)=\I\{L(x)\ge R(x)\}.
 \label{eq:ge-rule}
\end{equation}
The alternative rule is
\begin{equation}
 h_{>}(x)=\I\{L(x)>R(x)\}.
 \label{eq:gt-rule}
\end{equation}
Both are affine linear classifiers. If
\[
 w=(2^{m-1},\ldots,1,-2^{m-1},\ldots,-1)^\top,
\]
then
\begin{equation}
 h_{\ge}(x)=\I\{\langle w,x\rangle\ge0\},
 \qquad
 h_{>}(x)=\I\{\langle w,x\rangle-\tfrac12\ge0\}.
 \label{eq:linear-comparisons}
\end{equation}

Let \(\rho\) be uniform on \(\cX\). The labels are nearly balanced:
\begin{equation}
 \P(h_{\ge}(X)=1)=\frac12+2^{-m-1},
 \qquad
 \P(h_{>}(X)=1)=\frac12-2^{-m-1}.
 \label{eq:balanced-labels}
\end{equation}
The two rules disagree exactly on ties, so
\begin{equation}
 \P(h_{\ge}(X)\ne h_{>}(X))
 =\P(L(X)=R(X))=2^{-m}.
 \label{eq:tie-mass}
\end{equation}

To expose the learner symmetry, let \(s\) swap the first and second
\(m\)-bit blocks, and let \(c\) flip the binary label. Let \(g=(s,c)\) be the
task transformation with input component \(s\) and label component \(c\).
The induced action on rules gives
\begin{align}
 (gh_{\ge})(x)
 &=c\bigl(h_{\ge}(s^{-1}x)\bigr)\notag\\
 &=1-\I\{R(x)\ge L(x)\}
 =\I\{L(x)>R(x)\}
 =h_{>}(x).
 \label{eq:comparison-transform}
\end{align}
The uniform input law is invariant under the block swap. A learner whose
training rule is equivariant under the block swap and label flip therefore
treats \(h_{\ge}\) and \(h_{>}\) as orbit alternatives.

\begin{corollary}[Exponential passive-data requirement for one rule]
\label{cor:comparison-lower}
Let \(\cA\) be equivariant under the combined transformation \(g=(s,c)\). When
the inputs are independent and uniform on \(\{0,1\}^{2m}\),
\begin{equation}
 \P(H_{D_{h_{\ge}}}\ne h_{\ge})
 \ge\frac12(1-2^{-m})^n.
 \label{eq:comparison-failure}
\end{equation}
In particular, the failure probability exceeds \(1/4\) when
\(n<2^{m-1}\).
\end{corollary}
\begin{proof}
Equation \eqref{eq:comparison-transform} identifies the distinct orbit
alternative, and \cref{eq:tie-mass} gives its disagreement mass. Apply
\cref{thm:target-symmetry}.
\end{proof}

The corollary applies to a concrete gradient learner. Consider a fully
connected scalar-score network trained with binary cross-entropy by gradient
descent or minibatch SGD. Suppose that the first-layer initialization is
exchangeable across input coordinates, the output-layer initialization is
invariant under sign reversal, and regularization, optimizer state, stopping,
preprocessing, and hard-decision ties preserve the combined action. The
parameter map that swaps the two blocks of first-layer columns and negates the
output layer is a signed coordinate permutation. The calculation in
\cref{sec:gradient-symmetry} proves equivariance under \(g=(s,c)\), so
\cref{eq:comparison-failure} is a lower bound for this gradient-trained
network on the fixed target \(h_{\ge}\).

Returning \(h_{>}\) instead of \(h_{\ge}\) has statistical error only
\(2^{-m}\) under \(\rho\). The error is exponentially small even though the
returned rule is not the intended rule. Full support, balanced labels, and a
linear representation therefore do not resolve the exact-learning problem.

The designer can remove the information obstruction immediately. A single
chosen query with \(L(x)=R(x)\) receives label one under \(h_{\ge}\) and label
zero under \(h_{>}\). The exponential bound for uniform passive sampling does
not apply to this query protocol. A law concentrated on a tie has disagreement
mass one, so one observation misses the disagreement region with probability
zero. The example isolates the useful information supplied by a teacher: the
teacher knows which distinction the data must communicate.

The comparison convention can instead be encoded directly when the formal
operator \(\ge\) is already available to the system. Such an implementation
uses legitimate task knowledge and need not preserve the symmetry exchanging
\(\ge\) and \(>\). The learning problem is substantive when the system must
infer how an instruction, symbol, or natural-language request maps to the
comparison behavior. In that setting, a tie example communicates information
that the rest of the passive sample lacks.

\section{Target-aware teaching}
\label{sec:teaching}

The symmetry theorem is a design result as well as a lower bound. The theorem
locates the missing distinction, so the designer can decide where to supply it. The
distinction may enter through examples, queries, features, architecture,
initialization, a regularizer, or a verifier.

\Needspace{8\baselineskip}
\begin{definition}[Teaching set for a specified learner]
Fix a target rule \(h^\star\) and a deterministic learner \(\cA\). A labeled
dataset \(D\) is a \emph{teaching set} for \((\cA,h^\star)\) when
\begin{equation}
 \cA(D)=h^\star.
 \label{eq:teaching-set}
\end{equation}
\end{definition}

The definition depends on both objects. A dataset that teaches a
maximum-margin classifier need not teach a neural network trained by a
different objective. The teacher's knowledge of \(h^\star\) is also part of
the information protocol.

The following result gives a finite-dimensional positive example. The result
states the maximum-margin learner completely, including its treatment of the
intercept, so the teaching claim has a unique meaning.

Let \(\cX\subset\R^p\) be finite, and let
\(f:\cX\to\{-1,1\}\) be an affine-linearly separable target with both labels.
For a labeled dataset \(D\), write \(D_+\) and \(D_-\) for its positive and
negative input sets. Define
\begin{equation}
 w_D=\argmin_{w\in\R^p}\frac12\lVert w\rVert_2^2
 \quad\text{subject to}\quad
 \langle w,x^+-x^-\rangle\ge2
 \quad\text{for every }x^+\in D_+,\ x^-\in D_-.
 \label{eq:pairwise-margin}
\end{equation}
The strictly convex objective makes \(w_D\) unique. Set
\begin{equation}
 b_D=-\frac12\left(
     \min_{x^+\in D_+}\langle w_D,x^+\rangle
    +\max_{x^-\in D_-}\langle w_D,x^-\rangle
   \right),
 \label{eq:margin-intercept}
\end{equation}
and return \(f_D(x)=1\) when
\(\langle w_D,x\rangle+b_D\ge0\), and \(-1\) otherwise.

\begin{proposition}[A small teaching set for the maximum-margin learner]
\label{prop:max-margin-teaching}
Let \(D_0=\{(x,f(x)):x\in\cX\}\). There is a subset
\(D\subseteq D_0\) with
\begin{equation}
 |D|\le2p+2
 \label{eq:teaching-size}
\end{equation}
such that
\begin{equation}
 w_D=w_{D_0},\qquad b_D=b_{D_0},\qquad
 f_D(x)=f(x)\quad\text{for every }x\in\cX.
 \label{eq:teaching-conclusion}
\end{equation}
\end{proposition}

The proof in \cref{app:max-margin} selects at most \(p+1\) active
positive--negative pairs that certify the full-data maximum-margin direction.
Each pair contributes at most two examples. The result contrasts with passive
sampling: target-aware selection can communicate the intended rule using a
small subset of the domain even when random data rarely reveal the decisive
distinctions.

The proposition does not assert that a small teaching set is easy to find for
an arbitrary trained system. Computing a teaching set may require solving a
difficult design problem, and a more flexible learner can have more
symmetries or a less transparent selection rule. The proposition establishes
the possible benefit of matching the data to the target and the learner.

\subsection{Where task knowledge can enter}
\label{sec:task-knowledge}

The comparison-rule example makes four interventions concrete.

\begin{enumerate}[leftmargin=*]
\item \textbf{Target-aware data.} Include equality cases or other inputs from
the relevant disagreement region.
\item \textbf{Active queries.} Ask for a label or counterexample at an input
chosen to separate the current rule from the intended rule.
\item \textbf{Task-specific structure.} Encode the roles of the two operands,
the comparison convention, or other known asymmetries in the representation or
learning procedure.
\item \textbf{Verification.} Search for an input on which a proposed rule
violates the specification, then return that input as training information.
\end{enumerate}

The four interventions answer the course's recurring source-of-competence
question. The missing information can be supplied by the designer, the data
selection procedure, interaction, or a verifier. Statistical training alone
does not create a distinction that the complete training setup leaves
symmetric.

\section{Exact components and statistical interfaces}
\label{sec:systems}

Exact behavior of a complete AI system need not come entirely from one
exactly learned predictor. A statistical natural-language interface may map a
request to a formal representation, after which a sound executor or verifier
performs an exact operation.

Suppose that a parser produces a formal object \(\widehat S\) from an input
request whose correct formalization is \(S\). If an executor is correct
whenever \(\widehat S=S\), then
\begin{equation}
 \{\text{system output is incorrect}\}
 \subseteq
 \{\widehat S\ne S\}\cup
 \{\text{executor or invocation fails}\}.
 \label{eq:system-error-events}
\end{equation}
Consequently,
\begin{equation}
 \P(\text{system output is incorrect})
 \le
 \P(\widehat S\ne S)
 +\P(\text{executor or invocation fails}).
 \label{eq:system-error-bound}
\end{equation}
An exact executor removes one source of error. The complete system remains
statistical until semantic parsing, specification, tool selection, and
invocation are also controlled.

The symmetry theorem should therefore not be read as a general indictment of
statistical learning or as a representation impossibility. Neural networks
can represent many exact computations, as Lecture 14 showed. The theorem
identifies an information problem for a specified learner, target, and passive
data law. Task-specific training, interaction, and hybrid systems can change
that problem.

The next lecture studies denoising and score estimation. The denoising setting has a
different target: a population denoising objective identifies a conditional
mean and, through it, a distributional score. The relevant questions are
identifiability and estimation of a probabilistic object, followed by error
accumulation in reverse dynamics. Exact recovery of one deterministic rule is
not the universal criterion for learning.

\section{Optional reading: related symmetry claims}
\label{sec:optional}

Two related lines of work illustrate why claim scope matters. Empirical work
on recursive feature machines reports that particular algebraic train--test
partitions preserve symmetries that prevent generalization, while breaking
those data symmetries changes the observed behavior. These experiments concern
specified kernels, tasks, and training procedures; they are not the theorem
proved in \cref{thm:target-symmetry}.

Language-generation impossibility results instead quantify over collections
of possible languages and specified positive-example protocols. Such results
can show that one procedure cannot satisfy consistency and breadth throughout
a language family. They do not by themselves establish that a development
process using knowledge of English cannot build an English-speaking system.
The target-specific symmetry theorem supplies the missing kind of connection:
it states an obstruction for one target only after the target, learner, data
law, and relevant symmetry have been identified.

\section{Take-home messages}
\label{sec:take-home}

\begin{enumerate}[leftmargin=*]
\item Small population risk under one input law and correctness on every valid
input are different guarantees. Exactness for one specified task does not
require one learning procedure to succeed for every possible target task.
\item Pairwise indistinguishability yields a hypothesis-class lower bound:
some member of a close pair must be hard. The class result alone does not
lower-bound failure on the one intended target, because a learner can favor
that target and sacrifice the other member of the pair.
\item A justified learner symmetry prevents a representation or implementation
convention from deciding a question that should be settled by data or an
explicit task specification. Label names, record order, and spatial
transformations are symmetries only when they preserve the semantics,
observation protocol, loss, and available prior knowledge of the task.
\item Gradient descent inherits a learner symmetry when an orthogonal parameter
action preserves the predictor and training objective and the initialization
law is invariant. Optimizer state, preprocessing, stopping, and hard-decision
tie rules are also part of the symmetry claim.
\item The target-specific symmetry theorem lower-bounds exact recovery for one
fixed target, learner, and data law. The target orbit is a proof object induced
by the learner's symmetry. Symmetry turns the pairwise sum bound into an equal
failure probability for each orbit-related target.
\item For the rule \(L\ge R\), block-swap and label-flip symmetry produces the
alternative \(L>R\). Their disagreement mass is \(2^{-m}\), so a symmetric
learner needs exponentially many passive uniform examples to recover the fixed
intended rule with high probability.
\item Target-aware examples, active queries, task-specific structure, and
verification can break the relevant symmetry. A small teaching set is
learner-specific: its effectiveness depends on both the intended rule and the
selection rule used by the learner. A desirable learner symmetry can remain in
place while the teaching protocol supplies symmetry-breaking evidence.
\item Exact executors and verifiers can supply exact components inside a larger
statistical system. Their guarantees do not automatically control errors in
semantic parsing, specification, or tool invocation.
\end{enumerate}

\printendnotes

\section*{Bibliographical notes}

\paragraph{Statistical and exact learning.}
The distinction between average statistical performance and exact correctness,
the comparison-rule example, the generic disagreement lower bound, the
target-specific symmetry theorem, and the maximum-margin teaching result are
based on \citet{gyorgy2025}. The paper states its symmetry result using
nonidentity group elements. \Cref{eq:target-disagreement} instead ranges over
elements satisfying \(gh^\star\ne h^\star\), because a nonidentity element in
the target's stabilizer creates no alternative rule. The proof is otherwise the
same. The exact bound \cref{eq:target-symmetry-bound} keeps the missed-region
probability before the paper simplifies it to the constant \(1/4\).

The optional calculation in \cref{app:kl-directions} specializes the
conditional-support discussion from Lecture 4. The calculation is retained as
a comparison, not as part of the live argument. Finite reverse KL to a
deterministic full-support target characterizes exact conditional support, but
does not provide graded feedback among imperfect candidates. A cost-defined
Gibbs target provides graded feedback only when an evaluable cost supplies the
missing information.

\paragraph{Teaching and active learning.}
Classical exact-learning models use membership and equivalence queries;
\citet{angluin1988} is a foundational reference. Teaching instead permits
target-aware selection of examples for a specified learner. The
maximum-margin proof in \cref{app:max-margin} makes the learner's direction and
intercept choices explicit; it is a type-complete formulation of the support
constraint argument in \citet{gyorgy2025}. \citet{liu2016} study teaching
dimensions for several linear learners.

\paragraph{Implicit bias.}
For linearly separable data, gradient descent on logistic or cross-entropy loss
approaches the maximum-margin direction under the assumptions studied by
\citet{soudry2018}. The connection does not make every training set a teaching
set: the maximum-margin solution must itself implement the intended rule on
the entire valid domain. Optimization can also approach that direction slowly,
so repairing the information problem does not automatically give a fast
parameter optimization procedure.

\paragraph{Gradient-learner symmetry.}
The parameter-coupling argument in \cref{prop:gd-equivariance} supplies a
direct test for symmetry of a gradient learner. The claims about multilayer
perceptrons, binary heads, and vocabulary renamings require all of the listed
conditions; a model architecture alone does not establish learner symmetry.
\citet{gyorgy2025} use input-variable, label, and token symmetries to connect
the target-specific theorem to gradient-trained multilayer perceptrons and
transformers. The present notes state the parameter action and optimizer
qualifications explicitly.

\paragraph{Related symmetry results.}
\citet{tomasbernal2026} report an empirical symmetry-breaking mechanism for
recursive feature machines on algebraic tasks. The paper proves algebraic facts
about the task symmetry groups and reports the learning behavior empirically.
\citet{kalavasis2025} study consistency and breadth in language generation
under language-family and information-protocol assumptions. The optional
discussion uses both works only within those stated scopes.

\section*{Glossary}

\begin{description}[style=nextline,leftmargin=1.5em,labelindent=0pt]
\item[Statistical error \(L_\rho(h)\)]
The probability that \(h\) disagrees with the fixed target \(h^\star\) on an
input drawn from the specified law \(\rho\).
\item[Exact correctness]
Agreement with the fixed intended rule on every valid input in \(\cX\).
\item[Pairwise indistinguishability]
If a passive sample misses the disagreement region of two rules, the labeled
data are identical under both rules. The sum of their exact-recovery failure
probabilities is therefore at least the probability of missing that region.
\item[Worst-target class failure \(\Phi_n(\cA,\rho,\cH)\)]
The supremum of the exact-recovery failure probability over targets in
\(\cH\). A close pair lower-bounds this quantity without identifying which
member of the pair is hard.
\item[Group action on rules]
The transformation \((gh)(x)=g(h(g^{-1}x))\), which changes a rule consistently
with the group's actions on inputs and labels.
\item[Invariant input law]
A law \(\rho\) for which transforming an input by any \(g\in G\) preserves its
distribution.
\item[Equivariant learner]
A learning procedure whose output on transformed data has the same law as the
correspondingly transformed output on the original data.
\item[Parameter action \(S_g\)]
An orthogonal transformation of the parameter vector that realizes a data
transformation at the predictor level and preserves the training objective.
\item[Target orbit \(\cO_G(h^\star)\)]
The set \(\{gh^\star:g\in G\}\) of transformed alternatives induced by the
learner's symmetry.
\item[Target stabilizer]
The transformations \(g\in G\) satisfying \(gh^\star=h^\star\). Such a
transformation does not create a competing rule.
\item[Orbit disagreement \(\alpha_g(h^\star)\)]
The probability under \(\rho\) that the target and its transformed alternative
\(gh^\star\) give different labels.
\item[Target-specific disagreement \(\Delta_{\rho,G}(h^\star)\)]
The infimum orbit disagreement over transformations that produce a rule
different from the target.
\item[Teaching set]
A target-dependent labeled dataset that causes a specified learner to return
the intended rule.
\item[Maximum-margin learner]
The learner defined in \cref{eq:pairwise-margin,eq:margin-intercept}; it selects
the minimum-norm separating direction and a specified midpoint intercept.
\end{description}

\section{Exercises}
\label{sec:exercises}

\exerciseinstructions
\setcounter{courseexercise}{0}
\begin{exercisechunk}
\item \textbf{Hard and graded reverse KL.}\suppmark
\label[exercise]{ex:l22-reverse-kl-exactness}
\exerciseprofile{2}{2}{3}{\exproof\quad\excalculation\quad\exinterpretation}{%
Identify why hard reverse KL supplies no finite graded comparison among
imperfect rules and what additional information makes one possible.}
Let $\mathcal U$ be finite, let $\mu(u)>0$ for every $u\in\mathcal U$, and
let $f:\mathcal U\to\{0,1\}$. Define
\[
 P_f(u,y)=\mu(u)\mathbb I\{y=f(u)\},
 \qquad
 Q_q(u,y)=\mu(u)q(y\mid u),
\]
where $q(\,\cdot\mid u)$ is a conditional mass function.
\begin{enumerate}[(a)]
\item Prove that $D(Q_q,P_f)<+\infty$ if and only if
$q(f(u)\mid u)=1$ for every $u\in\mathcal U$. Determine the value of the
divergence when it is finite.
\item Prove
\[
 D(P_f,Q_q)=\sum_{u\in\mathcal U}\mu(u)
 \log\frac{1}{q(f(u)\mid u)}.
\]
Fix $u_0\in\mathcal U$ and $0<\eta<1$. Let
$q(1-f(u_0)\mid u_0)=\eta$ and let $q$ return the correct answer with
probability one at every other prompt. Compute both KL divergences.
\item For fixed $\eta$, explain how the forward divergence can be made
arbitrarily small by changing only $\mu(u_0)$. Explain why the same
construction never satisfies the reverse-KL support requirement.
\item Fix $0<\epsilon<1/2$ and define
\[
 P_{\epsilon,f}(u,y)
 =\mu(u)\bigl((1-\epsilon)\mathbb I\{y=f(u)\}
              +\epsilon\mathbb I\{y=1-f(u)\}\bigr).
\]
For $e_q(u)=q(1-f(u)\mid u)$, prove
\[
 D(Q_q,P_{\epsilon,f})
 =\sum_{u\in\mathcal U}\mu(u)\left(
 (1-e_q(u))\log\frac{1-e_q(u)}{1-\epsilon}
 +e_q(u)\log\frac{e_q(u)}{\epsilon}\right).
\]
Show that the contribution at $u$ diverges as $\epsilon\downarrow0$ when
$e_q(u)>0$. Explain why smoothing makes the objective finite and graded but
does not reveal $f(u)$ at an unobserved prompt.
\item Let $r(\,\cdot\mid u)$ be a positive reference mass function, let
$c:\mathcal U\times\{0,1\}\to\mathbb R$ be an evaluable cost, and let
$\beta>0$. Define
\[
 P_{\beta,c}(u,y)
 =\mu(u)\frac{r(y\mid u)e^{-\beta c(u,y)}}{Z_\beta(u)},
 \qquad
 Z_\beta(u)=\sum_y r(y\mid u)e^{-\beta c(u,y)}.
\]
Prove
\[
 D(Q_q,P_{\beta,c})
 =\sum_{u\in\mathcal U}\mu(u)\left(
 D(q(\,\cdot\mid u),r(\,\cdot\mid u))
 +\beta\sum_yq(y\mid u)c(u,y)+\log Z_\beta(u)\right).
\]
Now take $c(u,y)=\mathbb I\{y\ne f(u)\}$. Show that the minimizing
conditional mass function assigns positive probability to the incorrect
answer at every input for every finite $\beta$. Conclude that the graded Gibbs
objective is not by itself an exact-learning guarantee.

Identify the information, beyond finite positive demonstrations, that permits
an optimization method to evaluate the $q$-dependent terms on generated
prompt--answer pairs.
\end{enumerate}
\begin{hints}
A KL term is infinite when its first argument assigns positive mass to an
event to which its second argument assigns zero mass. For the two graded
targets, substitute their mass functions into the definition of KL and cancel
the common prompt mass $\mu(u)$.
\end{hints}
\end{exercisechunk}

\begin{exercisechunk}
\item \textbf{Symmetries of a gradient-trained linear classifier.}
\label[exercise]{ex:l22-gradient-symmetry}
\exerciseprofile{2}{2}{3}{\exproof\quad\excalculation\quad\exinterpretation}{%
Verify a symmetry of the complete training procedure, including the loss,
initialization, regularization, and hard-decision rule.}
Let
\[
 z_{w,b}(x)=\langle w,x\rangle+b,
 \qquad
 \ell_{\mathrm{ce}}(z,y)=\log(1+\exp(z))-yz,
\]
where $x,w\in\mathbb R^d$ and $y\in\{0,1\}$. For a labeled dataset
$D=((x_i,y_i))_{i=1}^n$, define
\[
 F_D(w,b)=\frac1n\sum_{i=1}^n
 \ell_{\mathrm{ce}}(z_{w,b}(x_i),y_i)
 +\frac\lambda2\lVert w\rVert_2^2.
\]
Let $P$ be a permutation matrix. Write $PD=((Px_i,y_i))_{i=1}^n$ and
$\bar D=((x_i,1-y_i))_{i=1}^n$.
\begin{enumerate}[(a)]
\item Prove
\[
 F_{PD}(Pw,b)=F_D(w,b),
 \qquad
 F_{\bar D}(-w,-b)=F_D(w,b).
\]
Deduce the corresponding transformations of the gradients.
\item Apply full-batch gradient descent with a fixed stepsize sequence.
Prove that runs initialized at $(w_0,b_0)$, $(Pw_0,b_0)$, and
$(-w_0,-b_0)$ on $D$, $PD$, and $\bar D$, respectively, remain related by
the same transformations at every iteration.
\item Suppose that $w_0$ has independent identically distributed centered
Gaussian coordinates and that $b_0$ has a centered Gaussian law independent
of $w_0$. Explain why the randomized learner has input-coordinate and
binary-label symmetry. Which property of the initialization law is used for
each symmetry?
\item Replace the isotropic penalty by
$\tfrac12\sum_{j=1}^d\lambda_jw_j^2$. Give a necessary and sufficient
condition on $(\lambda_1,\ldots,\lambda_d)$ for the particular permutation
$P$ to remain a symmetry. Explain why the deterministic hard decision
$\mathbb I\{z\ge0\}$ breaks label symmetry at $z=0$, and give an equivariant
tie rule.
\end{enumerate}
\begin{hints}
Use $\ell_{\mathrm{ce}}(-z,1-y)=\ell_{\mathrm{ce}}(z,y)$. Differentiate an
objective identity only after accounting for the orthogonal parameter map.
\end{hints}
\end{exercisechunk}

\begin{exercisechunk}
\item \textbf{A fixed comparison rule and a symmetric learner.}
\label[exercise]{ex:l22-comparison-symmetry}
\exerciseprofile{3}{2}{3}{\exproof\quad\excalculation\quad\exinterpretation}{%
Apply learner symmetry to a specified target rather than to a target class.}
Fix $m\ge1$. For $x\in\{0,1\}^{2m}$, let $L(x)$ and $R(x)$ be the integers
encoded by the first and second blocks of $m$ bits. Let
\[
 h_{\ge}(x)=\mathbb I\{L(x)\ge R(x)\},
 \qquad
 h_>(x)=\mathbb I\{L(x)>R(x)\}.
\]
Let $X$ be uniform on $\{0,1\}^{2m}$. Let $s$ swap the two bit blocks, let
$c$ flip a binary label, and let $g=cs$ act on a rule by
$(gh)(x)=c(h(s^{-1}x))$.
A randomized learner $\mathcal A$ is \emph{equivariant under $g$} if, for
every labeled dataset $D$, the output $\mathcal A(gD)$ has the same law as
$g\mathcal A(D)$, where $gD$ transforms every input by $s$ and every label
by $c$.
\begin{enumerate}[(a)]
\item Prove that $gh_{\ge}=h_>$. Determine
$\mathbb P(h_{\ge}(X)\ne h_>(X))$ and the positive-label probability of each
rule.
\item Let $\mathcal A$ be equivariant under $g$. Prove directly, using a
coupling of the two labeled-sample experiments, that the exact-recovery
failure probabilities of $\mathcal A$ on $h_{\ge}$ and $h_>$ are equal.
\item Prove that, when $\mathcal A$ is trained on $n$ independent uniform
examples labeled by the fixed intended rule $h_{\ge}$,
\[
 \mathbb P(\mathcal A\text{ does not return }h_{\ge})
 \ge \frac12(1-2^{-m})^n.
\]
Deduce that the failure probability is greater than $1/4$ when
$n<2^{m-1}$.
\item Give one target-aware membership query whose answer distinguishes the
two rules. Identify the assumption of the passive-data lower bound that the
query changes.
\end{enumerate}
\begin{hints}
The two rules disagree exactly when the encoded integers are equal. For part
(b), transform every input by $s$ and every label by $c$.
\end{hints}
\end{exercisechunk}

\begin{exercisechunk}
\item \textbf{A small teaching set for a maximum-margin learner.}\suppmark
\label[exercise]{ex:l22-max-margin-teaching}
\exerciseprofile{3}{3}{3}{\exproof\quad\exconstruction\quad\exinterpretation}{%
Use active constraints to teach a specified target to a specified learner.}
Let $\mathcal X\subset\mathbb R^p$ be finite, and let
$f:\mathcal X\to\{-1,1\}$ be affine-separable with both labels present.
The full labeled dataset is
$D_0=\{(x,f(x)):x\in\mathcal X\}$. For a labeled subset $D$, define
$D_+=\{x:(x,1)\in D\}$ and $D_-=\{x:(x,-1)\in D\}$. The learner chooses
\[
 w_D=\arg\min_w\frac12\lVert w\rVert_2^2
 \quad\text{subject to}\quad
 \langle w,x^+-x^-\rangle\ge2
 \quad\text{for every }x^+\in D_+,\ x^-\in D_-,
\]
and then sets
\[
 b_D=-\frac12\left(
 \min_{x^+\in D_+}\langle w_D,x^+\rangle
 +\max_{x^-\in D_-}\langle w_D,x^-\rangle\right).
\]
The returned rule is $x\mapsto\operatorname{sign}(\langle w_D,x\rangle+b_D)$.
You may use the conic Carath\'eodory theorem: every vector in a cone generated
in $\mathbb R^p$ can be expressed as a nonnegative combination of at most
$p$ generators. Using the weaker bound $p+1$ is also acceptable.
\begin{enumerate}[(a)]
\item Let $w_0=w_{D_0}$. Use the first-order optimality condition to show
that $w_0$ is a nonnegative combination of at most $p+1$ active difference
vectors $x_j^+-x_j^-$.
\item Let $D$ contain the endpoints of the selected active differences.
Prove $|D|\le2p+2$ and $w_D=w_0$.
\item Prove that the selected positive endpoints attain
$\min_{x^+\in(D_0)_+}\langle w_0,x^+\rangle$ and that the selected negative
endpoints attain
$\max_{x^-\in(D_0)_-}\langle w_0,x^-\rangle$. Deduce $b_D=b_{D_0}$ and
prove that the learner returns $f$ on all of $\mathcal X$.
\item Explain why the construction does not provide one universal dataset
that teaches every affine rule or every learning algorithm.
\end{enumerate}
\begin{hints}
The normal cone of the feasible polyhedron at $w_0$ is generated by the
negatives of its active constraint normals. For (c), pair a hypothetical
smaller positive projection with a selected negative endpoint.
\end{hints}
\end{exercisechunk}


\appendix
\section{Optional comparison: KL direction and exact support}
\label{app:kl-directions}

The main lecture uses disagreement events because they expose the information
missing from passive data. KL divergence gives a related support calculation,
but the calculation does not produce a finite-data exact-learning method.

Let \(\cU\) be finite, let \(\mu(u)>0\) for every \(u\in\cU\), and let
\(f:\cU\to\{0,1\}\) be a deterministic answer rule. Define
\begin{equation}
 P_f(u,y)=\mu(u)\I\{y=f(u)\},
 \qquad
 Q_q(u,y)=\mu(u)q(y\mid u),
 \label{eq:appendix-kl-laws}
\end{equation}
where \(q(\,\cdot\mid u)\) is a conditional mass function. The two KL
directions answer different questions:
\begin{align}
 \KL(Q_q,P_f)<+\infty
 &\quad\Longleftrightarrow\quad
 q(f(u)\mid u)=1
 \quad\text{for every }u\in\cU,
 \label{eq:appendix-reverse-exact}\\
 \KL(P_f,Q_q)
 &=\sum_{u\in\cU}\mu(u)\log\frac{1}{q(f(u)\mid u)}.
 \label{eq:appendix-forward-kl}
\end{align}
Indeed, positive candidate mass on \((u,1-f(u))\) creates an infinite
reverse-KL term because \(P_f(u,1-f(u))=0\). The forward direction instead
averages the log loss of the correct response under \(\mu\). A rare input can
therefore contribute little to forward KL while still preventing exact
correctness. Hard reverse KL detects the failure, but assigns the same infinite
value to every candidate that assigns any positive mass to an incorrect
response.

An empirical target does not resolve the missing-information problem. If
\(\widehat P_n\) is the empirical law of \(n\) demonstrations, then
\begin{equation}
 \KL(Q,\widehat P_n)<+\infty
 \quad\Longrightarrow\quad
 \operatorname{supp}(Q)
 \subseteq\operatorname{supp}(\widehat P_n).
 \label{eq:appendix-empirical-support}
\end{equation}
The right side contains only observed prompt--answer pairs. The empirical
reverse-KL objective therefore supplies no intended response at an unseen
prompt.

Smoothing the target support makes reverse KL finite and graded. The smoothed
law still depends on \(f(u)\) at every input, including the inputs absent from
the demonstrations. Smoothing changes the penalty assigned to an error; it
does not reveal the intended answer. Part (d) of
\cref{ex:l22-reverse-kl-exactness} gives the corresponding binary-KL formula.

A cost-defined Gibbs target gives a useful graded objective when a method can
evaluate generated responses. Let \(r(\,\cdot\mid u)\) be positive, let
\(c(u,y)\) be an evaluable cost, and let \(\beta>0\). Define
\begin{equation}
 P_{\beta,c}(u,y)
 =\mu(u)\frac{r(y\mid u)e^{-\beta c(u,y)}}{Z_\beta(u)},
 \qquad
 Z_\beta(u)=\sum_y r(y\mid u)e^{-\beta c(u,y)}.
 \label{eq:appendix-gibbs-target}
\end{equation}
Then
\begin{equation}
 \KL(Q_q,P_{\beta,c})
 =\sum_{u\in\cU}\mu(u)\left(
 \KL(q(\,\cdot\mid u),r(\,\cdot\mid u))
 +\beta\sum_y q(y\mid u)c(u,y)+\log Z_\beta(u)\right).
 \label{eq:appendix-gibbs-identity}
\end{equation}
After division by \(\beta\), minimizing the right side is equivalent to
minimizing expected generated cost plus a
\(\beta^{-1}\)-scaled KL penalty relative to \(r\). The evaluable cost, not the
choice of KL direction, supplies feedback on generated answers.

The graded Gibbs objective also does not give exactness at a finite value of
\(\beta\). For the correctness cost
\(c_f(u,y)=\I\{y\ne f(u)\}\), the minimizing conditional law is the conditional
part of \(P_{\beta,c_f}\), and
\begin{equation}
 q_\beta^\star(1-f(u)\mid u)
 =\frac{r(1-f(u)\mid u)e^{-\beta}}
 {r(f(u)\mid u)+r(1-f(u)\mid u)e^{-\beta}}>0
 \label{eq:appendix-gibbs-not-exact}
\end{equation}
for every finite \(\beta\). Taking \(\beta\) large can make the incorrect mass
small, but evaluating \(c_f\) already requires the ability to check the
intended answer. The Gibbs construction is useful for KL-regularized
expected-cost optimization when such an evaluator exists. The construction is
not a solution to exact learning from finite demonstrations alone.

\section{Proof of the maximum-margin teaching result}
\label{app:max-margin}

The appendix proves the positive teaching result from
\cref{sec:teaching}. The proof selects a small set of active constraints that
certifies the same strictly convex optimization solution as the full labeled
domain.

\begin{proposition}[A small teaching set for the maximum-margin learner]
\label{prop:max-margin-teaching-restated}
Under the assumptions and learner definition in
\cref{eq:pairwise-margin,eq:margin-intercept}, let
\(D_0=\{(x,f(x)):x\in\cX\}\). There is a subset
\(D\subseteq D_0\) with \(|D|\le2p+2\) such that
\[
 w_D=w_{D_0},\qquad b_D=b_{D_0},\qquad
 f_D(x)=f(x)\quad\text{for every }x\in\cX.
\]
\end{proposition}
\begin{proof}
For a dataset \(S\) containing both labels, define the positive--negative
difference set
\[
 \Gamma(S)=\{x^+-x^-:x^+\in S_+,\ x^-\in S_-\}
\]
and the feasible polyhedron
\[
 C(S)=\{w\in\R^p:\langle w,v\rangle\ge2
                    \text{ for every }v\in\Gamma(S)\}.
\]
Write \(w_0=w_{D_0}\). First-order optimality for minimizing
\(\lVert w\rVert_2^2/2\) over \(C(D_0)\) gives
\begin{equation}
 -w_0\in N_{C(D_0)}(w_0),
 \label{eq:normal-cone-condition}
\end{equation}
where \(N_C(w)\) is the normal cone of \(C\) at \(w\).

Only active constraints generate this normal cone. Thus
\begin{equation}
 -w_0\in\operatorname{cone}
 \{-v:v\in\Gamma(D_0),\ \langle w_0,v\rangle=2\}.
 \label{eq:active-cone}
\end{equation}
By the conic form of Carath\'eodory's theorem, one may select at most
\(p+1\) active differences \(v_1,\ldots,v_k\), \(k\le p+1\), whose negative
cone still contains \(-w_0\). For each \(v_j\), choose a pair
\((x_j^+,x_j^-)\) satisfying \(v_j=x_j^+-x_j^-\), and let \(D\) contain
these labeled endpoints. Then \(|D|\le2p+2\).

The selected constraints imply
\(-w_0\in N_{C(D)}(w_0)\). Therefore \(w_0\) satisfies the first-order
optimality condition for the reduced problem. Strict convexity gives
\begin{equation}
 w_D=w_0.
 \label{eq:same-direction}
\end{equation}

The intercept remains to be checked. Every selected difference is active for
the full problem:
\[
 \langle w_0,x_j^+\rangle-\langle w_0,x_j^-\rangle=2.
\]
All full-data positive--negative differences have value at least two. Hence
each selected positive endpoint attains
\(\min_{x^+\in(D_0)_+}\langle w_0,x^+\rangle\), and each selected negative
endpoint attains
\(\max_{x^-\in(D_0)_-}\langle w_0,x^-\rangle\). Otherwise pairing a smaller
full-data positive point or a larger full-data negative point with the other
selected endpoint would violate feasibility. Consequently the extrema used in
\cref{eq:margin-intercept} are unchanged by restricting from \(D_0\) to
\(D\), so \(b_D=b_{D_0}\).

The full-data classifier labels every point in \(\cX\) according to \(f\).
The equal direction and intercept therefore give \(f_D=f\) on \(\cX\).
\end{proof}

\clearpage
\begin{thebibliography}{99}

\bibitem[Angluin(1988)]{angluin1988}
Dana Angluin.
\newblock Queries and concept learning.
\newblock \emph{Machine Learning}, 2(4):319--342, 1988.
\newblock \url{https://doi.org/10.1007/BF00116828}.

\bibitem[Gy\"orgy et~al.(2025)Gy\"orgy, Lattimore, Lazi\'c, and
Szepesv\'ari]{gyorgy2025}
Andr\'as Gy\"orgy, Tor Lattimore, Nevena Lazi\'c, and Csaba Szepesv\'ari.
\newblock Beyond statistical learning: Exact learning is essential for
general intelligence.
\newblock arXiv:2506.23908, 2025.
\newblock \url{https://arxiv.org/abs/2506.23908}.

\bibitem[Kalavasis et~al.(2025)Kalavasis, Mehrotra, and
Velegkas]{kalavasis2025}
Alkis Kalavasis, Anay Mehrotra, and Grigoris Velegkas.
\newblock On the limits of language generation: Trade-offs between
hallucination and mode collapse.
\newblock In \emph{Proceedings of the 57th Annual ACM Symposium on Theory of
Computing}, pages 1732--1743, 2025.
\newblock \url{https://doi.org/10.1145/3717823.3718108}.

\bibitem[Liu and Zhu(2016)]{liu2016}
Ji Liu and Xiaojin Zhu.
\newblock The teaching dimension of linear learners.
\newblock \emph{Journal of Machine Learning Research}, 17(162):1--25, 2016.
\newblock \url{https://jmlr.org/papers/v17/15-153.html}.

\bibitem[Soudry et~al.(2018)Soudry, Hoffer, Nacson, Gunasekar, and
Srebro]{soudry2018}
Daniel Soudry, Elad Hoffer, Mor Shpigel Nacson, Suriya Gunasekar, and
Nathan Srebro.
\newblock The implicit bias of gradient descent on separable data.
\newblock \emph{Journal of Machine Learning Research}, 19(70):1--57, 2018.
\newblock \url{https://jmlr.org/papers/v19/18-188.html}.

\bibitem[Tom\`as Bernal et~al.(2026)Tom\`as Bernal, Mallinar, and
Belkin]{tomasbernal2026}
Marcel Tom\`as Bernal, Neil Rohit Mallinar, and Mikhail Belkin.
\newblock Breaking data symmetry is needed for generalization in feature
learning kernels.
\newblock In \emph{Proceedings of the 29th International Conference on
Artificial Intelligence and Statistics}, PMLR 300:4204--4212, 2026.
\newblock \url{https://proceedings.mlr.press/v300/bernal26a.html}.

\end{thebibliography}

\end{document}
